\documentclass[../../hsm_tok_q.tex]{subfiles}

\begin{document}
    \setcounter{figure}{0}
    Sさんのグループは，先生が示した問題をみんなで考えた．

    次の各問に答えよ．

    \begin{tcolorbox}[title=［先生が示した問題］,
        colback=white, coltitle=black, colbacktitle=white,colframe=black, 
        toptitle=1pt, bottomtitle=1pt,
        arc=4pt, boxrule=0.5pt, before upper={\parindent=1\zw},
        left=3pt, right=3pt, top=3pt, bottom=4pt]

    \textsf{図\ref{fig/hsm_tok_2024_1st_02_01_q}}のように，円Oの円周を12等分する点に，
    1から12までの自然数の番号を，小さい順で時計回りに付ける．

    1から12までの番号を付けた点のうち，2点を結んでできる線分が円Oの直径となるとき，
    その2点を向かい合う点とする．

    例えば，1の点と7の点は向かい合う点である．

    図\textsf{\ref{fig/hsm_tok_2024_1st_02_01_q}}において，1組の向かい合う点を選び，それぞれの点の番号のうち，
    小さい方の数を$a$，大きい方の数を$b$とする．

    $a$，$b$の平均値をA，$b^2-a^2$の値をBとするとき，BはAの何倍か求めなさい．
    %
    \vspace{.5\baselineskip}
    {\par \centering
        \includegraphics[width=.65\linewidth]%
            {\subfix{fig_hsm_tok_2024_1st_02/fig_hsm_tok_2024_1st_02_01_q.pdf}}
            \captionof{figure}{} \label{fig/hsm_tok_2024_1st_02_01_q}
    \par}
    %
    \end{tcolorbox}
    \begin{enumerate}
        \item ［先生が示した問題］で，BはAの\:\myboxa{　　　}\:倍と表すとき，
            \:\myboxa{　　　}\:に当てはまる数を，次の{\textsf ア}〜\textsf{エ}のうちから
            選び，記号で答えよ．
    
    \begin{tabularx}{\dimexpr\linewidth-3.5mm}{@{}LLLL@{}}
        {\textsf ア}\quad$-3$ & {\textsf イ}\quad$4$ & {\textsf ウ}\quad$6$ & {\textsf エ}\quad$12$
    \end{tabularx}
    \end{enumerate}

    \newpage
    
    Sさんのグループは，［先生が示した問題］をもとにして，次の問題を作った．

    \begin{tcolorbox}[title=［Sさんのグループが作った問題］,
        colback=white, coltitle=black, colbacktitle=white,colframe=black, 
        toptitle=1pt, bottomtitle=1pt,
        arc=4pt, boxrule=0.5pt, before upper={\parindent=1\zw},
        left=3pt, right=3pt, top=3pt, bottom=4pt]

    \textsf{\ref{fig/hsm_tok_2024_1st_02_02_q}}のように，円Oの円周を24等分する点に，
    1から24までの自然数の番号を，小さい順で時計回りに付ける．

    1から24までの番号を付けた点のうち，2点を結んでできる線分が円Oの直径となるとき，
    その2点を向かい合う点とする．

    図\textsf{\ref{fig/hsm_tok_2024_1st_02_02_q}}において，
    異なる2組の向かい合う点を選び，1組目のそれぞれの点の番号のうち，
    小さい方の数を$a$，大きい方の数を$b$とし，
    2組目のそれぞれの点の番号のうち，小さい方の数を$c$，大きい方の数を$d$とする．
    %
    \vspace{.5\baselineskip}
    {\par \centering
        \includegraphics[width=.65\linewidth]%
            {\subfix{fig_hsm_tok_2024_1st_02/fig_hsm_tok_2024_1st_02_02_q.pdf}}
            \captionof{figure}{} \label{fig/hsm_tok_2024_1st_02_02_q}
    \par}
    %
    \end{tcolorbox}
    %
    \begin{enumerate}[resume]
        \item $a$，$b$，$c$，$d$の平均値をP，$bd-ac$の値をQとするとき，
            $\mathrm{Q}=24\mathrm{P}$となることを証明せよ．
    \end{enumerate}
\end{document}

